\begin{derivation}[有限脉冲散射的侧带振幅]
固定一个电子切片 $(\zeta,\tau)$，把 $\beta(\zeta,\tau)$ 暂记为 $\beta$，并令
\[
 \theta\equiv\kappa_\omega\zeta.
\]
局域散射算符是纯相位
\begin{equation*}
 U(\theta)=\exp\!\left(\beta \ee^{\ii\theta}-\beta^*\ee^{-\ii\theta}\right).
\end{equation*}
由于两个指数中的量只是复数，它们彼此对易，所以可精确分解为
\begin{equation*}
 U(\theta)
 =\ee^{\beta \ee^{\ii\theta}}
  \ee^{-\beta^*\ee^{-\ii\theta}}.
\end{equation*}
分别作 Taylor 展开，得到双重级数
\begin{align*}
 U(\theta)
 &=\left[\sum_{r=0}^{\infty}
 \frac{\beta^r}{r!}\ee^{\ii r\theta}\right]
 \left[\sum_{s=0}^{\infty}
 \frac{(-\beta^*)^s}{s!}\ee^{-\ii s\theta}\right]\\
 &=\sum_{r,s\ge0}
 \frac{\beta^r(-\beta^*)^s}{r!s!}
 \ee^{\ii(r-s)\theta}.
\end{align*}
这里 $r$ 与 $s$ 分别计数从两个指数中取出的正、负频因子数；真正由能量测量分辨的是它们的差
\[
 \ell=r-s,
\]
而不是 $r$、$s$ 各自。于是同一个最终边带 $\ell$ 包含无穷多条具有相同净交换量、但往返次数不同的振幅路径。

若 $\ell\ge0$，所有非负整数解为
\[
 r=j+\ell,\qquad s=j,
 \qquad j=0,1,2,\ldots;
\]
若 $\ell<0$，则
\[
 r=j,\qquad s=j-\ell.
\]
引入
\[
 \ell_+=\max(\ell,0),
 \qquad
 \ell_-=\max(-\ell,0),
\]
两种情况统一为
\[
 r=j+\ell_+,
 \qquad
 s=j+\ell_-.
\]
因此第 $j$ 条路径对净边带 $\ell$ 的复振幅贡献为
\begin{equation*}
 C_j^{(\ell)}
 =\frac{\beta^{j+\ell_+}(-\beta^*)^{j+\ell_-}}
 {(j+\ell_+)!(j+\ell_-)!},
\end{equation*}
恢复 $(\zeta,\tau)$ 依赖即得到正文\eqnrefs{eq:park_Cjn_general}。这里 $C_j^{(\ell)}$ 是振幅而不是概率；所有具有相同 $\ell$ 的路径必须先相干相加：
\begin{equation*}
 A_\ell(\zeta,\tau)
 =\sum_{j=0}^{\infty}C_j^{(\ell)}(\zeta,\tau).
\end{equation*}

这一求和可以显式完成。写
\[
 \beta=|\beta|\ee^{\ii\varphi_\beta}.
\]
对 $\ell\ge0$，有
\begin{align*}
 A_\ell
 &=\beta^\ell
 \sum_{j=0}^{\infty}
 \frac{(-|\beta|^2)^j}{j!(j+\ell)!}\\
 &=\ee^{\ii\ell\varphi_\beta}
 \sum_{j=0}^{\infty}
 \frac{(-1)^j|\beta|^{2j+\ell}}{j!(j+\ell)!}\\
 &=\ee^{\ii\ell\varphi_\beta}J_\ell(2|\beta|),
\end{align*}
其中最后一步使用 Bessel 函数的级数定义
\[
 J_\ell(2x)
 =\sum_{j=0}^{\infty}
 \frac{(-1)^j x^{2j+\ell}}{j!(j+\ell)!}.
\]
对 $\ell<0$ 使用 $J_{-n}(x)=(-1)^nJ_n(x)$ 得到同一统一结果
\begin{equation*}
 \boxed{
 A_\ell=\ee^{\ii\ell\varphi_\beta}J_\ell(2|\beta|).}
\end{equation*}
因此 Bessel 边带并不是额外假设，而是同一净交换通道内所有往返路径振幅的精确重求和。

还可以直接检验概率守恒。由于指数
$\beta \ee^{\ii\theta}-\beta^*\ee^{-\ii\theta}$
纯虚，任意 $\theta$ 都有 $|U(\theta)|=1$。Fourier 系数定义为
\begin{equation*}
 A_\ell=\frac{1}{2\pi}
 \int_0^{2\pi}\dd\theta\,
 U(\theta)\ee^{-\ii\ell\theta}.
\end{equation*}
Parseval 恒等式于是给
\begin{align*}
 \sum_{\ell\in\mathbb Z}|A_\ell|^2
 &=\frac{1}{2\pi}\int_0^{2\pi}\dd\theta\,|U(\theta)|^2\\
 &=1.
\end{align*}
\end{derivation}

\begin{derivation}[快分支振荡与长作用时间条件]
把归一化开关包络取在有限区间 $[t_i,t_f]$，满足
\[
 \widetilde w(0)=\int_{t_i}^{t_f}\dd t\,w(t)=1,
 \qquad
 \widetilde w(\Omega)=\int_{t_i}^{t_f}\dd t\,w(t)\ee^{\ii\Omega t}.
\]
正文的一阶微扰已经把慢支路与快支路之比化为
$|\widetilde w(2\omega)|$。为了看清“作用时间长”为什么能够压低它，假设 $w$ 在作用区间绝对连续。对 $\Omega\ne0$ 分部积分：
\begin{align*}
 \widetilde w(\Omega)
 &=\left.\frac{w(t)\ee^{\ii\Omega t}}{\ii\Omega}\right|_{t_i}^{t_f}
 -\frac1{\ii\Omega}\int_{t_i}^{t_f}\dd t\,
 \dot w(t)\ee^{\ii\Omega t}.
\end{align*}
取模并用三角不等式得到严格上界
\begin{align*}
 |\widetilde w(\Omega)|
 \le\frac{|w(t_f)|+|w(t_i)|
 +\int_{t_i}^{t_f}\dd t\,|\dot w(t)|}{|\Omega|}.
\end{align*}
因此在 $\Omega=2\omega$ 处，若包络端点平滑降到零，便有
\begin{equation*}
 \left|\frac{\mathcal A_{\rm fast}^{(1)}}{\mathcal A_{\rm slow}^{(1)}}\right|
 =|\widetilde w(2\omega)|
 \le\frac{1}{2\omega}
 \int_{t_i}^{t_f}\dd t\,|\dot w(t)|.
\end{equation*}
这个式子比单纯说“$T$ 很长”更精确：快支路由包络在光学频率尺度上的变化率控制；同样持续时间的两个脉冲，只要边缘平滑程度不同，反旋泄漏就可以不同。

矩形窗口不满足端点为零的平滑条件，需要直接积分。对
$w(t)=T^{-1}\mathbf 1_{[-T/2,T/2]}(t)$，
\begin{align*}
 \widetilde w(2\omega)
 &=\frac1T\int_{-T/2}^{T/2}\dd t\,\ee^{2\ii\omega t}\\
 &=\frac{\sin(\omega T)}{\omega T}
 =\sinc(\omega T).
\end{align*}
所以正文\eqnrefs{eq:feist2015_long_interaction_condition} 的一般小参数是
$\epsilon_{2\omega}=|\widetilde w(2\omega)|$；只有指定矩形开关之后，它才退化为由 $\omega T$ 单独控制的 sinc 条件。
\end{derivation}
